\documentclass[11pt,a4paper]{article}
\usepackage[utf8]{inputenc}
\usepackage[T1]{fontenc}
\usepackage[french]{babel}
\usepackage{geometry}
\geometry{margin=2cm}
\usepackage{xcolor}
\usepackage{hyperref}
\usepackage{graphicx}
\usepackage{listings}
\usepackage{tabularx}
\usepackage{booktabs}
\usepackage{amsmath}

\definecolor{primary}{RGB}{37, 99, 235}
\definecolor{dark}{RGB}{15, 23, 42}
\definecolor{lightbg}{RGB}{248, 250, 252}

\hypersetup{
    colorlinks=true,
    linkcolor=primary,
    urlcolor=primary
}

\lstset{
    basicstyle=\ttfamily\small,
    backgroundcolor=\color{lightbg},
    frame=single,
    breaklines=true
}

\title{\textbf{\Huge 🛡️ GuardRail DevSecOps}\\[0.3em] \Large Rapport Technique d'Ingénierie — Phases 1 \& 2\\Durcissement Sécurité \& Cycle de Vie Persistant des Vulnérabilités}
\date{4 Septembre 2026}

\begin{document}

\maketitle

\begin{abstract}
Ce rapport détaille les réalisations techniques des Phases 1 et 2 sur la plateforme DevSecOps \textbf{GuardRail} (Rails 8 API-only, PostgreSQL, SolidQueue, Brakeman, Semgrep). L'intégralité du socle d'authentification, de contrôle d'accès multi-tenant (IDOR/BOLA), de limitation de débit (Rack::Attack), de calcul d'empreinte stable SHA-256 et du cycle de vie persistant des vulnérabilités a été validée avec un taux de réussite de \textbf{100\% (138 tests passés, 0 échec)}.
\end{abstract}

\section{Phase 1 — Durcissement de la Sécurité}

\subsection{Isolation Multi-Tenant \& Prévention IDOR (GR-101 / GR-102)}
L'audit complet des 12 contrôleurs a confirmé l'absence de contournement d'authentification (\texttt{User.first}, \texttt{current\_user ||}). Une vulnérabilité critique d'injection de clé étrangère a été neutralisée dans \texttt{RepositoriesController} via le filtre \texttt{validate\_project\_ownership}, interdisant l'association d'un dépôt à un projet d'un autre utilisateur.

\subsection{Limitation de Débit \& Anti-DDoS (GR-103)}
Intégration de \texttt{rack-attack} avec seuils stricts renvoyant un code HTTP 429 et l'entête standard \texttt{Retry-After} :
\begin{itemize}
    \item \texttt{POST /api/v1/login} : 5 req / 20s (anti brute-force)
    \item \texttt{POST /api/v1/signup} : 3 req / 60s (anti création de comptes)
    \item \texttt{POST /api/v1/webhooks/github} : 30 req / 60s (anti-flood)
    \item \texttt{ALL /api/*} : 100 req / 60s (anti-DDoS global)
\end{itemize}

\subsection{Correction de la Suite de Tests \& Tests de Sécurité (GR-104)}
Résolution de 16 régressions préexistantes et ajout de specs d'isolation validant le rejet systématique (HTTP 422) des tentatives d'injection de clé étrangère.

\subsection{Pagination SQL Native (GR-105)}
Module \texttt{Paginatable} validé sur toutes les collections avec requêtes \texttt{LIMIT}/\texttt{OFFSET} et métadonnées complètes.

\section{Phase 2 — Persistance \& Cycle de Vie des Vulnérabilités}

\subsection{Algorithme d'Empreinte Stable (GR-201)}
L'empreinte ne dépend plus du numéro de ligne absolu, évitant les fausses alertes lorsqu'un développeur ajoute du code :
\begin{equation}
\text{Fingerprint} = \text{SHA256}(\text{scanner} \parallel \text{"|"} \parallel \text{rule\_id} \parallel \text{"|"} \parallel \text{file} \parallel \text{"|"} \parallel \text{normalized\_code})
\end{equation}

\subsection{Cycle de Vie Immuable \& Détection des Réouvertures (GR-202 / GR-203)}
Suppression totale de \texttt{destroy\_all} (0 occurrence dans le code). Les scans constituent des instantanés historiques immuables.
\begin{itemize}
    \item \textbf{OPEN} : Détectée pour la première fois.
    \item \textbf{IGNORED} : Faux positif / risque accepté reporté automatiquement sans jamais être rouvert.
    \item \textbf{RESOLVED} : Corrigée et archivée dans l'historique sans destruction.
    \item \textbf{REOPENED} : Détectée à nouveau après avoir été résolue, avec horodatage \texttt{reopened\_at} et réactivation des pénalités de score.
\end{itemize}

\section{Résultats des Tests \& Métriques}
\begin{itemize}
    \item \textbf{RSpec} : 138 exemples exécutés, \textbf{0 échec}, 1 pending.
    \item \textbf{RuboCop} : Code 100\% conforme aux standards Rails.
    \item \textbf{Base de Données} : Contrainte d'unicité \texttt{(scan\_id, fingerprint)} validée.
\end{itemize}

\section{Prochaine Étape Recommandée (Phase 3)}
Pipeline Multi-Scanner \& SCA : intégration de \texttt{bundler-audit}, \texttt{npm-audit} et \texttt{Gitleaks}.

\end{document}
